\documentclass[conference]{IEEEtran}
\IEEEoverridecommandlockouts

\usepackage{cite}
\usepackage{amsmath,amssymb,amsfonts}
\usepackage{algorithmic}
\usepackage{graphicx}
\usepackage{textcomp}
\usepackage{xcolor}
\usepackage{booktabs}
\usepackage{multirow}
\usepackage{url}
\usepackage{hyperref}
\usepackage{microtype}

\def\BibTeX{{\rm B\kern-.05em{\sc i\kern-.025em b}\kern-.08em
    T\kern-.1667em\lower.7ex\hbox{E}\kern-.125emX}}

\begin{document}

\title{Learning Compiler Pass Sequences with GRU-Guided Reinforcement Learning}

\author{
\IEEEauthorblockN{[Author Name]}
\IEEEauthorblockA{\textit{Department of Computer Science} \\
[University Name] \\
[Email]}
}

\maketitle

% ─────────────────────────────────────────────────────────────────────────────
\begin{abstract}
Selecting an effective sequence of LLVM optimization passes is a combinatorial
problem that conventional fixed pipelines (e.g., \texttt{-O3}) solve only
approximately.  We present an early-stage knowledge discovery system that
learns pass-ordering policies from empirical execution-time measurements of 38
hand-written C benchmarks.  Starting from a corpus of 950\,000 randomly sampled
pass sequences (up to 20 passes from a curated menu of 28 primary transforms),
we perform exploratory data analysis to characterize the optimization landscape,
quantify per-pass impact, and identify benchmarks whose performance is most
sensitive to pass ordering.  Preliminary random-search results show that 18 of
38 benchmarks already achieve lower execution times than \texttt{-O3} without
any learned policy, with speedups of $-27\%$ on \texttt{interpreter} and
$-24\%$ on \texttt{kmp\_search}.  The full system will train a gated recurrent
unit (GRU) policy network with proximal policy optimization (PPO) on a curated
subset of six benchmarks, targeting pass-sequence policies that generalize
across unseen programs.  This midterm report describes the data collection
infrastructure, preprocessing choices, exploratory analysis, and a project
plan for the remaining training and evaluation phases.
\end{abstract}

\begin{IEEEkeywords}
compiler optimization, phase ordering, reinforcement learning, GRU, LLVM,
knowledge discovery
\end{IEEEkeywords}

% ─────────────────────────────────────────────────────────────────────────────
\section{Introduction}

\subsection{Background}

Modern optimizing compilers such as LLVM~\cite{lattner2004llvm} expose
hundreds of individual transformation passes---inlining, loop unrolling,
scalar replacement of aggregates, vectorization, dead-code elimination, and
many others.  A \emph{pass pipeline} is an ordered sequence of these
transformations applied to a program's intermediate representation (IR).
Order matters: inlining before constant propagation exposes more constants to
fold; loop rotation must precede loop-invariant code motion (LICM) for LICM to
fire on the canonical loop form.

LLVM's \texttt{-O3} flag encodes a fixed pipeline tuned by compiler engineers
for average-case performance across a broad program distribution.  While
effective in aggregate, this one-size-fits-all schedule is suboptimal for
individual programs: the same benchmark compiled with a different pass ordering
of equal length can run substantially faster or slower.  Finding the best
pipeline for a specific program is NP-hard in general~\cite{almagor2004finding}
and has been studied as the \emph{phase-ordering problem}.

Machine learning offers a principled alternative: learn a policy that maps
program features to pass choices, trained on real execution-time feedback.  If
the policy captures which program characteristics predict the benefit of each
pass, it can generalize to unseen programs and outperform hand-crafted
schedules.

\subsection{Related Work}

Early work applied iterative compilation and machine learning to phase
ordering~\cite{fursin2011milepost,ashouri2017micomp}.  MILEPOST
GCC~\cite{fursin2011milepost} learned a nearest-neighbor predictor from static
program features to select optimization flags, achieving consistent speedups
over \texttt{-O3} on MiBench.

More recently, deep reinforcement learning has produced compelling results.
AutoPhase~\cite{huang2019autophase} formulated phase ordering as a Markov
decision process (MDP) and trained an LSTM policy with PPO on LLVM IR features,
demonstrating speedups over \texttt{-O3} on the LLVM test suite.
CompilerGym~\cite{cummins2021compilergym} provided a standardized OpenAI Gym
interface to compiler optimization environments, enabling reproducible RL
research.  MLGO~\cite{trofin2021mlgo} demonstrated that a learned inlining
policy deployed inside production LLVM matches or exceeds the hand-tuned
heuristic at Google scale.

Our work differs from AutoPhase in two key respects.  First, we target
wall-clock execution time measured with hardware timing (\texttt{CLOCK\_MONOTONIC})
rather than proxy metrics such as IR instruction count, producing noisier but
more faithful reward signals.  Second, we collect a large offline corpus of
random sequences prior to training, enabling systematic exploratory data
analysis to characterize the landscape and guide architecture choices before
any learning is attempted.  We also adopt a GRU rather than an LSTM, reducing
parameter count and training time while retaining sequential memory of past
pass decisions.

\subsection{Dataset Description}

The dataset consists of 38 self-contained C benchmark programs spanning
diverse computational patterns: sorting and searching
(\texttt{quicksort}, \texttt{binary\_search}, \texttt{kmp\_search}), numerical
computation (\texttt{fft}, \texttt{polynomial\_eval},
\texttt{matrix\_multiply\_tiled}), data-structure operations
(\texttt{binary\_tree}, \texttt{heap\_ops}, \texttt{hashtable}), string and
compression codecs (\texttt{rle\_codec}, \texttt{lz\_compress},
\texttt{huffman}), and interpreted or recursive workloads
(\texttt{interpreter}, \texttt{tail\_recursive}, \texttt{nqueens}).

Each benchmark is a single C file that performs its computation inside a
\texttt{run()} function called repeatedly by a shared timing harness.  All
benchmarks share \texttt{bench\_timing.h}, which implements a 10\%-trimmed
mean of 201 timed iterations (preceded by 5 untimed warm-up executions) using
\texttt{CLOCK\_MONOTONIC}.  Trimming the extreme deciles discards outliers
caused by OS page faults, TLB misses, and CPU frequency transients that would
otherwise inflate variance.  The harness prints the trimmed-mean execution
time in nanoseconds to standard output.

For each benchmark and each randomly sampled pass sequence, the compilation
pipeline executes four steps:
\begin{enumerate}
  \item \textbf{IR emission}: compile the C source to annotated IR with
        \texttt{clang-20 -O3 -Xclang -disable-llvm-optzns}.  The \texttt{-O3}
        frontend flag activates Clang's analysis infrastructure (type-based
        alias analysis metadata, lifetime markers, loop hint annotations),
        while \texttt{-disable-llvm-optzns} suppresses all LLVM optimization
        passes.  The result is an unoptimized IR enriched with analysis
        annotations that subsequent passes can exploit.
  \item \textbf{Pass application}: apply the sampled sequence with
        \texttt{opt-20 -passes=\textit{pipeline}}.
  \item \textbf{Native compilation}: lower the optimized IR to a binary with
        \texttt{clang-20 -O3 -Xclang -disable-llvm-passes}, activating O3
        backend code generation (instruction selection, register allocation,
        scheduling) without re-running any IR optimization passes.
  \item \textbf{Benchmarking}: execute the binary and record the trimmed-mean
        time printed to stdout.
\end{enumerate}
Passes are always applied from the same annotated base IR to ensure
determinism; no incremental IR state is carried between trials.

The corpus comprises 950\,000 records: 25\,000 random sequences per benchmark,
each sequence drawn uniformly from the 28-pass primary menu with length sampled
uniformly from 1 to 20.  Three baseline compilations
(\texttt{-O0}, \texttt{-O2}, \texttt{-O3}) are also recorded for each
benchmark.

\subsection{Model Design}

The system is an actor--critic RL agent with a GRU backbone
(Figure~\ref{fig:arch}).  At each step $t$ the agent observes a state
$\mathbf{s}_t = (\mathbf{f}, \mathbf{h}_{t-1})$ where
$\mathbf{f} \in \mathbb{R}^{18}$ is a static feature vector extracted from the
current program IR and $\mathbf{h}_{t-1}$ is the GRU hidden state encoding
the history of previously applied passes.  The policy network outputs a
categorical distribution over 29 actions (28 passes $+$ STOP); the value
network outputs a scalar baseline for variance reduction.  Training uses
PPO~\cite{schulman2017proximal} with a clipped surrogate objective.

The GRU architecture is motivated by the sequential, order-dependent nature of
pass pipelines.  A stateless (feed-forward) policy cannot distinguish
\emph{inline before sroa} from \emph{sroa before inline}; the GRU hidden state
provides a compact memory of past decisions.  GRU was preferred over LSTM for
this task because the gating mechanism is simpler (two gates versus three),
reducing parameter count and forward-pass latency---important when each
environment step incurs a compilation and benchmarking round that already
dominates wall-clock time.  Prior work~\cite{huang2019autophase} confirmed that
recurrent architectures outperform feed-forward baselines on phase ordering.

\begin{figure}[t]
  \centering
  \fbox{\parbox{0.9\linewidth}{\centering\vspace{1em}
    \textbf{GRU Actor--Critic Architecture} \\[0.6em]
    IR features $(18\text{-d})$ $\xrightarrow{\text{Linear}}$ input dim $H$ \\
    Action embedding $(d\text{-d})$ $\rightarrow$ concatenate \\[0.3em]
    $\downarrow$ GRU layer(s), hidden size $H$ $\downarrow$ \\[0.3em]
    Policy head $\rightarrow$ 29 action logits \\
    Value head $\rightarrow$ scalar $V(\mathbf{s})$
    \vspace{1em}
  }}
  \caption{GRU actor--critic architecture.  IR features are projected to the
    GRU input dimension and concatenated with a learned embedding of the
    previous action.  The GRU hidden state carries the pass history.  The
    policy head selects the next pass; the value head estimates expected return
    for PPO variance reduction.}
  \label{fig:arch}
\end{figure}

% ─────────────────────────────────────────────────────────────────────────────
\section{Data Preprocessing and Exploratory Data Analysis}

\subsection{Pass Menu Curation}

A na\"ive action space containing all LLVM passes would be impractical: the
combinatorial explosion makes exploration intractable; many passes are
language-specific (e.g., coroutine transforms, OpenMP lowering) or
module-level analytics irrelevant to single-function C compute code; and some
passes are dominated variants of others already in the menu.

We defined a \emph{primary pass menu} of 28 transforms satisfying three
criteria simultaneously: (1) the pass appears inside LLVM's \texttt{-O3}
inner optimization loop or is a direct prerequisite; (2) it has demonstrable,
broad impact on typical C compute code; and (3) it covers a distinct
optimization category with no redundant substitute already in the menu.  The
28 passes span instruction cleanup (\texttt{instcombine}, \texttt{adce},
\texttt{dse}, \texttt{sccp}, \texttt{gvn}), memory and SSA promotion
(\texttt{mem2reg}, \texttt{sroa}, \texttt{memcpyopt}), CFG simplification
(\texttt{simplifycfg}, \texttt{jump-threading}), inlining (\texttt{inline}),
loop infrastructure and canonicalization (\texttt{loop-rotate}, \texttt{licm},
\texttt{indvars}, \texttt{loop-idiom}, \texttt{loop-deletion},
\texttt{loop-unswitch}), loop transformations (\texttt{loop-unroll},
\texttt{loop-vectorize}), and whole-function passes
(\texttt{tailcallelim}, \texttt{slp-vectorizer}, \texttt{early-cse}).
Several passes are included in two parameterization variants (e.g.,
\texttt{sroa} vs.\ \texttt{sroa<modify-cfg>},
\texttt{licm} vs.\ \texttt{licm<allowspeculation>}) so the policy can learn
when the more aggressive form is beneficial.

LLVM 20 imposes nesting constraints that required additional preprocessing:
\texttt{inline} must be wrapped as \texttt{cgscc(inline)};
\texttt{licm} must appear inside \texttt{loop-mssa(\ldots)}; and
\texttt{loop-rotate} is a LoopPass that must be wrapped in
\texttt{loop(\ldots)} to avoid consuming subsequent passes as loop-level
passes.  These nesting rules are handled transparently by the pass menu module
before the pipeline string is passed to \texttt{opt-20}.

\subsection{IR Feature Extraction}

For each compiled IR file the system extracts 18 features by a single-pass
text scan of the \texttt{.ll} file in under 50\,ms.  Features fall into two
groups:

\textbf{Instruction-type counts:} arithmetic (\texttt{add}, \texttt{mul}),
memory (\texttt{load}, \texttt{store}, \texttt{alloca}, \texttt{gep}), control
flow (\texttt{br}, \texttt{phi}), function calls (\texttt{call}), comparisons
(\texttt{icmp}, \texttt{fcmp}), returns, and a catch-all \emph{other}.

\textbf{Structural features:} basic-block count (including the implicit entry
block per function), total instruction count, function count, approximate loop
depth (estimated from back-edge detection in label order), and the
load-to-store ratio.

Because the base IR is emitted with \texttt{-O3 -disable-llvm-optzns}, the IR
carries rich analysis metadata from Clang's frontend: type-based alias analysis
(TBAA) annotations on loads and stores, \texttt{llvm.lifetime} intrinsics
marking stack allocation lifetimes, and loop hint metadata.  While we do not
currently use this metadata as features, it is available to the \texttt{opt-20}
passes applied during each episode, meaning passes such as \texttt{licm} and
\texttt{loop-vectorize} benefit from richer alias information than they would
receive from a bare \texttt{-O0} IR.

\subsection{Baseline Landscape}

Table~\ref{tab:baselines} summarizes the optimization headroom for the 38
benchmarks, sorted by the \texttt{-O3}/\texttt{-O0} speedup ratio.
\texttt{select\_chain} offers the largest headroom ($19.6\times$ speedup from
\texttt{-O0} to \texttt{-O3}), while \texttt{binary\_tree} and
\texttt{graph\_bfs} have only $1.6\times$--$1.7\times$, reflecting
pointer-chasing workloads where memory latency dominates over compiler
transformations.  Notably, \texttt{matrix\_multiply\_tiled} shows a larger
\texttt{-O2} speedup than \texttt{-O3} ($4.2\times$ vs.\ $6.2\times$ relative
to \texttt{-O0}), indicating that \texttt{-O3}'s additional vectorization
passes actively hurt this benchmark --- a result that motivates learning
program-specific sequences rather than applying the same aggressive pipeline
everywhere.

\begin{table}[t]
\caption{Baseline execution times and optimization headroom for representative
benchmarks. Times are 10\%-trimmed-mean nanoseconds.}
\label{tab:baselines}
\centering
\begin{tabular}{lrrrc}
\toprule
Benchmark & O0 (ns) & O2 (ns) & O3 (ns) & O3/O0 \\
\midrule
select\_chain      &   754\,252 &   38\,265 &   38\,416 & 19.6$\times$ \\
convolution        &   647\,015 &   46\,919 &   53\,609 & 12.1$\times$ \\
stencil2d          &   926\,164 &   81\,357 &   81\,900 & 11.3$\times$ \\
tail\_recursive    & 6\,119\,777 &  562\,730 &  573\,013 & 10.7$\times$ \\
dead\_store\_elim  & 1\,282\,547 &  129\,625 &  130\,070 &  9.9$\times$ \\
hashtable          &   405\,484 &   43\,761 &   41\,858 &  9.7$\times$ \\
matrix\_chain      &    29\,690 &    4\,099 &    4\,081 &  7.3$\times$ \\
deep\_call\_chain  &   534\,090 &   73\,569 &   73\,503 &  7.3$\times$ \\
\midrule
matrix\_mult.\ (tiled) & 648\,044 & 153\,817 & 105\,055 & 6.2$\times$ \\
\midrule
quicksort          & 1\,528\,784 &  767\,406 &  787\,649 &  1.9$\times$ \\
graph\_bfs         & 1\,244\,459 &  702\,294 &  714\,793 &  1.7$\times$ \\
binary\_tree       &   193\,579 &  118\,839 &  117\,990 &  1.6$\times$ \\
\bottomrule
\end{tabular}
\end{table}

\subsection{Performance Distributions and Coefficient of Variation}

Figure~\ref{fig:distributions} shows the distribution of execution times
across all 25\,000 random sequences per benchmark.  The coefficient of
variation (CV) measures how sensitive a benchmark is to pass ordering.  High
CV indicates a rugged optimization landscape and, consequently, a richer reward
signal for learning.

\begin{figure}[t]
  \centering
  \includegraphics[width=\linewidth]{../eda_output/distributions.png}
  \caption{Execution-time distributions for all 38 benchmarks across 25\,000
    random pass sequences.  Box plots show P25/P75 with whiskers at P10/P90.
    Benchmarks are sorted by coefficient of variation (CV).}
  \label{fig:distributions}
\end{figure}

\texttt{tail\_recursive} has the highest CV at 77.9\%, spanning from
1.06\,$\mu$s (P10) to 10.74\,$\mu$s (P90)---a $10\times$ range from worst to
best random sequence.  Other high-CV benchmarks include
\texttt{matrix\_multiply\_tiled} (62.2\%), \texttt{struct\_pack} (47.0\%),
\texttt{polynomial\_eval} (46.0\%), and \texttt{select\_chain} (46.7\%).
Low-CV benchmarks such as \texttt{interpreter} (13.9\%) and
\texttt{binary\_search} (15.1\%) are less sensitive to ordering---either
because a single dominant pass (\texttt{inline} for \texttt{interpreter})
decides most performance, or because the workload is memory-bound and largely
immune to instruction-level transformations.

\subsection{Pass Impact and Enrichment Analysis}

To quantify the marginal effect of each pass we compute two statistics:

\textbf{Geometric-mean speedup.}  For each pass $p$, we compare sequences
containing $p$ against sequences not containing $p$ (matched by length to
reduce confounding).  The geometric mean of the speedup ratios across all 38
benchmarks measures the unconditional benefit of including $p$.

\textbf{Top-decile enrichment.}  We identify the top-10\% of sequences by
execution time for each benchmark and compute the \emph{enrichment ratio}:
the presence rate of pass $p$ in top-10\% sequences divided by its overall
presence rate ($\approx 30\%$ under uniform random sampling with a 20-pass
length budget).

Figure~\ref{fig:enrichment} shows both metrics.  \texttt{inline} is the most
impactful pass by a wide margin: enrichment $2.35\times$, geo-speedup
$1.26\times$, present in 71\% of top-decile sequences versus 30\% overall.
This reflects inlining's role as a \emph{pass enabler}: eliminating call
overhead exposes callee bodies to the caller's context, allowing constant
propagation, dead-code elimination, and loop transformations to fire at a
wider scope.  \texttt{simplifycfg} ranks second ($1.68\times$ enrichment),
consistent with its role as a clean-up pass needed after nearly every other
transformation.  \texttt{sroa} and \texttt{mem2reg} rank third and fourth,
confirming that SSA promotion is a prerequisite for most downstream analyses.

\begin{figure}[t]
  \centering
  \includegraphics[width=\linewidth]{../eda_output/pass_enrichment.png}
  \caption{Pass enrichment (top-10\% presence rate vs.\ overall rate) for all
    28 primary passes.  All passes are positively enriched, confirming the
    primary menu is well-curated.}
  \label{fig:enrichment}
\end{figure}

\begin{figure}[t]
  \centering
  \includegraphics[width=\linewidth]{../eda_output/pass_impact.png}
  \caption{Geometric-mean speedup (across 38 benchmarks) for sequences
    containing each pass versus sequences not containing it.
    \texttt{inline} dominates; passes further right contribute incrementally.}
  \label{fig:impact}
\end{figure}

\subsection{Ceiling Analysis}

The ceiling analysis asks: what is the best execution time achievable within
the random-search budget, and how does it compare to \texttt{-O3}?
Figure~\ref{fig:ceiling} plots the gap between the best-found sequence and
\texttt{-O3} for all benchmarks.

\begin{figure}[t]
  \centering
  \includegraphics[width=\linewidth]{../eda_output/ceiling_gaps.png}
  \caption{Gap between best random sequence and \texttt{-O3} per benchmark.
    Negative values (blue) indicate random search beats \texttt{-O3}.
    Benchmarks in red require pipeline extensions not achievable with the
    current pass menu.}
  \label{fig:ceiling}
\end{figure}

Of 38 benchmarks, 18 are beaten by random search.  Four benchmarks
(\texttt{convolution}, \texttt{regex\_match}, \texttt{hashtable},
\texttt{matrix\_multiply\_tiled}) have gaps exceeding 25\%: these depend on
pass parameterizations not represented in our menu (e.g., loop-vectorize
interleave-width parameters or the four-iteration CGSCC devirtualization loop
used by \texttt{-O3}).  These four are excluded from the training set as the
agent cannot close the gap through pass ordering alone.

\subsection{Sequence Length Analysis}

Figure~\ref{fig:seqlength} shows the length of the best-found sequence for
each benchmark, grouped by difficulty tier.  The median best length is 17--18
across all tiers, and no spike occurs at the 20-pass cap, indicating that the
cap is not the binding constraint.  Several benchmarks achieve their best
results with substantially shorter sequences (e.g., \texttt{binary\_tree} at
length 9, \texttt{fft} at length 13), validating the design of a STOP action
that terminates an episode before the maximum.

\begin{figure}[t]
  \centering
  \includegraphics[width=\linewidth]{../eda_output/seq_length_by_tier.png}
  \caption{Distribution of best-found sequence lengths by difficulty tier.
    The 20-pass cap is not systematically binding; median best length is
    17--18.}
  \label{fig:seqlength}
\end{figure}

% ─────────────────────────────────────────────────────────────────────────────
\section{Preliminary Results}

\subsection{Random Search as Exploratory Mining}

Random pass-sequence search constitutes an unsupervised mining baseline: no
labels or human expertise guide the search, yet it discovers sequences that
outperform \texttt{-O3} on 18 of 38 benchmarks.  Table~\ref{tab:ceiling}
reports results for the six benchmarks selected for GRU training.

\begin{table}[t]
\caption{Best random-search result for the six training benchmarks.
  ``vs O3'' is the percentage improvement over \texttt{-O3} (negative = faster).
  CV is the coefficient of variation across all 25\,000 sequences.}
\label{tab:ceiling}
\centering
\begin{tabular}{lrrrc}
\toprule
Benchmark & O3 (ns) & Best (ns) & vs O3 & CV \\
\midrule
interpreter      &  37\,918 &  27\,591 & $-27.2\%$ & 13.9\% \\
kmp\_search      & 886\,408 & 671\,195 & $-24.3\%$ & 42.7\% \\
polynomial\_eval & 417\,235 & 364\,235 & $-12.7\%$ & 46.0\% \\
fft              &  46\,509 &  41\,852 & $-10.0\%$ & 35.5\% \\
array\_reduction &  97\,184 &  92\,619 &  $-4.7\%$ & 36.6\% \\
binary\_tree     & 117\,990 & 107\,087 &  $-9.2\%$ & 18.4\% \\
\bottomrule
\end{tabular}
\end{table}

The six benchmarks were selected to maximize diversity of optimization profile
(call-heavy, arithmetic, loop/vectorize, reduction, pointer/struct), coefficient
of variation (reward gradient for learning), and consistent ability to beat
\texttt{-O3} in random search, which confirms the pass menu is expressive
enough.

\subsection{Pass Presence in Best Sequences}

We examined which passes appear in the best-found sequences for each training
benchmark.  \texttt{inline} appears in all six; \texttt{simplifycfg} and
\texttt{sroa} appear in five of six.  For \texttt{interpreter} and
\texttt{kmp\_search}, \texttt{inline} appears as the first or second pass,
consistent with its role as an enabler.  For \texttt{fft} and
\texttt{array\_reduction}, \texttt{loop-vectorize} and \texttt{loop-unroll}
appear in all best sequences.

The absence of strong co-occurrence signals beyond individual enrichment (all
pair-lift values $< 1.1$) indicates that pass \emph{ordering} rather than pass
\emph{co-presence} is the primary lever.  This is precisely the structure a
GRU, with its sequential hidden state, is designed to capture.

% ─────────────────────────────────────────────────────────────────────────────
\section{Next Steps and Anticipated Results}

\subsection{GRU Policy Implementation}

The policy network will be implemented in Rust using the Burn deep learning
framework with the NDArray CPU backend (the model is small enough that GPU
transfer overhead outweighs computation savings).  The planned architecture is:

\begin{enumerate}
  \item \textbf{Input projection}: a linear layer mapping the 18 IR features
        (normalized to zero mean and unit variance per feature across the
        training corpus) to a hidden dimension $H \in \{64, 128\}$.
  \item \textbf{Action embedding}: a learned embedding of the previous action
        (pass index, dimensionality $d \in \{16, 32\}$) concatenated to the
        feature projection before the GRU input.
  \item \textbf{GRU}: one or two layers with hidden size $H$, processing the
        sequence of (feature projection $\oplus$ action embedding) vectors.
  \item \textbf{Policy head}: a linear layer from the GRU output to 29 action
        logits, with a temperature parameter and entropy bonus for exploration.
  \item \textbf{Value head}: a separate linear layer from the GRU output to a
        scalar return estimate.
\end{enumerate}

\subsection{Reward Function}

The reward signal assigns scaled points for achieving three performance tiers
relative to compiler baselines:

\begin{equation}
  r \;=\; w_0 \cdot \mathbf{1}[t < t_{\texttt{O0}}]
        + w_2 \cdot \mathbf{1}[t < t_{\texttt{O2}}]
        + w_3 \cdot \mathbf{1}[t < t_{\texttt{O3}}]
        + s_{\texttt{O3}} \cdot \Delta_{\texttt{O3}}
  \label{eq:reward}
\end{equation}

where $t$ is the agent's execution time, $\Delta_{\texttt{O3}} =
(t_{\texttt{O3}} - t)/t_{\texttt{O3}}$ is the fractional improvement over
\texttt{-O3} (positive when the agent wins), and $w_0, w_2, w_3, s_{\texttt{O3}}$
are tunable scaling weights.  The tiered structure ensures the agent receives
positive signal even when it only matches \texttt{-O0}, avoids the flat-zero
problem when \texttt{-O3} is not yet beaten, and provides a fine-grained
gradient for policies that surpass \texttt{-O3}.  All times are trimmed-mean
measurements from the same benchmark harness used during data collection.

\subsection{Training Protocol}

Training will follow the standard PPO recipe~\cite{schulman2017proximal} with
clipped surrogate objective ($\epsilon = 0.2$), generalized advantage
estimation ($\gamma = 0.99$, $\lambda = 0.95$), and an entropy bonus
($\beta = 0.01$) to encourage exploration.  The environment cycles over the
six selected benchmarks, sampling a new benchmark at the start of each
episode.

\subsection{Evaluation Plan}

The trained policy will be evaluated on: (1) the six training benchmarks
(in-distribution); (2) the remaining 28 benchmarks not seen during training
(out-of-distribution generalization); and (3) comparison baselines:
\texttt{-O0}, \texttt{-O2}, \texttt{-O3}, and the best random-search result
from the 25\,000-sequence corpus.  The primary metric is percentage improvement
over \texttt{-O3}; secondary metrics are the fraction of benchmarks that beat
\texttt{-O3} and the geometric-mean speedup over \texttt{-O3} across all 38
benchmarks.

\subsection{Anticipated Results}

Based on the EDA findings we anticipate:

\textbf{Training benchmarks.}  The agent should converge to within 5\% of the
random-search ceiling on all six benchmarks.  For \texttt{interpreter}
($-27\%$) and \texttt{kmp\_search} ($-24\%$), this would represent substantial
and reproducible improvements over \texttt{-O3}.

\textbf{Generalization.}  The enrichment analysis shows that a small set of
high-impact passes (\texttt{inline}, \texttt{simplifycfg}, \texttt{sroa},
\texttt{gvn}) benefit most programs.  A policy that learns to apply these
reliably early, then apply loop transforms when loop-heavy IR features are
present, should transfer to unseen benchmarks.  We expect the trained agent to
beat \texttt{-O3} on at least 20--25 of the 38 benchmarks.

\textbf{Failure modes.}  The four excluded benchmarks are expected to remain
below \texttt{-O3}; closing those gaps requires extending the pass menu with
additional parameterizations or reproducing the CGSCC devirtualization loop
structure used by \texttt{-O3}.

% ─────────────────────────────────────────────────────────────────────────────
\section*{Acknowledgment}
[Omitted for submission]

% ─────────────────────────────────────────────────────────────────────────────
\bibliographystyle{IEEEtran}
\bibliography{refs}

\end{document}
